% =====================================================================
\chapter{Lab 2 -- Sylvane Retraite: il canale di scambio file}
\label{lab:02}
% =====================================================================

\begin{nota}[Fonti e date]
Caso presentato da Gianluca il 12/10/2026 (Capitolo~\ref{lez:2026-10-12}), svolto venerdì 16/10/2026
(Capitolo~\ref{lez:2026-10-16}), discusso nel debrief di lunedì 19/10/2026. Materiale: README del caso,
\emph{Sylvane Retraite Breach Case Study} (PDF), \emph{Severity Assessment Worksheet} (obbligatorio),
\emph{Notification Form} (Excel, condizionale), immagine della \emph{curveball~1}. Caso di fantasia.
\textbf{Non esiste una soluzione ufficiale}: lo svolgimento ragionato è dei redattori ed è segnalato
come tale.
\end{nota}

\section{Il caso in sintesi}

\paragraph{Contesto.} Sylvane Retraite è un fondo pensione professionale francese (IRPS), titolare
dei dati dei propri iscritti. Un fornitore esterno (responsabile del trattamento) carica file
operativi su un \textbf{server SFTP nella DMZ del fondo}, che funge da ``cassetta delle lettere'':
il fondo sposta poi i file nei propri ambienti interni. L'accesso al server è con autenticazione a
\textbf{certificato} (niente password) e solo dall'IP pubblico del fornitore. Gli studenti sono il
\textbf{DPO}, che venerdì 11 aprile 2025 pomeriggio deve consigliare CEO, CISO e Head of Member
Services.

\paragraph{Cronologia (2025).}
\begin{center}
\small
\begin{tabularx}{\linewidth}{p{2.6cm} X}
\toprule
\textbf{Quando} & \textbf{Evento} \\ \midrule
Mer 2/4, 08:30 & Inizia l'accesso non autorizzato al server SFTP, attraverso i sistemi del fornitore e
  con le sue credenziali legittime. \\
2--4/4 & Circa 53 ore di permanenza; file esfiltrati. \\
Ven 4/4, 13:37 & Finestra chiusa: accesso bloccato, arrivano i forensi. Il fornitore viene a
  conoscenza dell'incidente e lo riconduce ai propri sistemi. \\
Sab 5 -- Gio 10/4 & ``Settimana silenziosa'': il fornitore completa i controlli; il fondo non sa nulla. \\
Ven 11/4, 10:44 & Comunicazione formale del fornitore via email certificata: ``qualcosa'' è stato
  esfiltrato, senza dettagli. Il fornitore ha già notificato la propria agenzia nazionale di
  cybersicurezza. \\
Ven 11/4, 11:12 & Il fondo classifica l'incidente come \textbf{grave ai sensi del DORA}; nessun
  impatto sulla continuità. \\
Ven 11/4, pomeriggio & Punto di decisione (Parte~A, ore 14:00). \\
\bottomrule
\end{tabularx}
\end{center}

\paragraph{Che cosa passava nel canale.}
\begin{center}
\small
\begin{tabularx}{\linewidth}{l X l}
\toprule
\textbf{Gruppo} & \textbf{Contenuto} & \textbf{Stato} \\ \midrule
G1 & nome e cognome dell'iscritto + valore totale della posizione pensionistica maturata (dati comuni +
  finanziari); oltre 100 iscritti, numero non ancora quantificato & confermato \\
G2 & file operativi caricati dal fornitore; categorie e volume da determinare & confermato \\
G3 & il server riceveva un \textbf{export giornaliero completo} del database: dati storici e recenti
  non separabili & indeterminato \\
G4 & quattro ambienti interni alimentati dalla DMZ (antiriciclaggio, BI, risk, middleware): non
  raggiunti & non coinvolto \\
G5 & categorie particolari (art.~9): nessuna individuata finora & nessuna \\
G6 & pubblicazione o abuso: nessuna evidenza (``l'assenza di prova non è prova di assenza'') & non
  accertato \\
\bottomrule
\end{tabularx}
\end{center}

\paragraph{Le tre posizioni.}
\begin{itemize}
  \item \textbf{CEO}: nessuna notifica, solo registro. L'accesso è passato dai sistemi del fornitore,
    il server è protetto da certificati, nulla è stato pubblicato; informare gli iscritti porterebbe
    solo un'ondata di telefonate. Si rivaluterà con nuove informazioni.
  \item \textbf{CISO}: solo CNIL. Nome, cognome e valore della posizione sono stati esfiltrati: c'è un
    rischio ex art.~33; notifica preliminare e integrazioni successive.
  \item \textbf{Head of Member Services}: CNIL e tutti gli iscritti coinvolti. Se un truffatore chiama
    un iscritto citando l'importo esatto della sua posizione (social engineering) e noi non abbiamo
    detto nulla, non abbiamo difese.
\end{itemize}

\section{Consegna}
\begin{itemize}
  \item \textbf{Tempi}: 105 minuti di \emph{wall clock} = tre giorni di \emph{exercise clock}. Le nuove
    informazioni (\emph{curveball}) arrivano in chat durante il lab.
  \item \textbf{Worksheet obbligatorio} per tutti, qualunque sia la conclusione (anche ``nessuna
    notifica''). Se la notifica è dovuta, si compila anche il \textbf{modulo di notifica} e si carica
    un unico zip con entrambi i file, non protetto da password.
  \item Nome del gruppo e dei componenti nell'intestazione e nel modulo di caricamento; file nominati
    con il nome del gruppo. Lavori non attribuibili: scartati.
  \item Compiti della Parte~A: (1) due perimetri (solo ciò che è stato preso o tutto ciò che il server
    conteneva); (2) punteggio dei gruppi di dati; (3) la severità non è il rischio; (4) scegliere chi
    informare: iscritti, CNIL, ACPR o nessuno; (5) da quando decorrono le 72 ore, e dalla conoscenza
    di chi.
  \item Indicazione data durante il lab: serve una valutazione \textbf{complessiva} della violazione,
    non un punteggio per ogni voce.
\end{itemize}

\subsection{Il worksheet}
Struttura (10 sezioni): sintesi dei fatti; DPC (tabella delle quattro categorie, con motivazione di
ogni scostamento); EI (0,25 / 0,50 / 0,75 / 1,00); CB (somma di riservatezza, integrità, disponibilità,
intento malevolo); calcolo di SE; fascia di rischio; decisione (solo registro / CNIL / CNIL e iscritti
/ segnalazione DORA all'ACPR); perimetro; \textbf{giustificazione} (la risposta da dare all'autorità
se chiede ``cosa avete fatto e perché''); \textbf{strategia} (il ragionamento interno, ``ciò che non
si può dire ad alta voce'').

\begin{nota}[Honne e tatemae]
Gianluca ha presentato le ultime due sezioni con due concetti giapponesi: \textbf{tatemae} è la
facciata, ciò che si mostra all'autorità (sezione~9); \textbf{honne} è la verità, il motivo reale
della scelta, che può includere anche il timore di essere licenziati (sezione~10). Nella realtà sono
due cose molto diverse.
\end{nota}

\subsection{Il modulo di notifica}
È la trasposizione, in Excel, del modulo online dell'autorità (art.~33): sezioni A (chi notifica: il
legale rappresentante o un delegato, con dichiarazione sulle conseguenze di dichiarazioni false),
B (notifica preliminare o completa), C (titolare), D (DPO), E (altri soggetti: il fornitore come
responsabile), F (informazioni sulla violazione: date, modalità di conoscenza, motivi del ritardo,
natura, causa, sistemi, misure in essere, categorie e numero di interessati e registrazioni),
H (misure adottate e per il futuro), L (comunicazione agli interessati, con il \textbf{testo effettivo}
della comunicazione), M (altre autorità: già indicata la notifica all'ACPR ai sensi del DORA),
N (trattamenti transfrontalieri). Alcuni campi sono precompilati perché lo studente non può
conoscerli.

\subsection{La curveball~1 (venerdì 11/4, 16:30)}
Il fornitore ha inviato una comunicazione identica a un'altra società del gruppo, che fa
\textbf{wealth management} per clienti facoltosi e di alto profilo. Stesso fornitore, stesso canale,
stessa finestra. Domande guida: il proprio perimetro e punteggio; chi è titolare dei dati della
società sorella; se il gruppo ha ora un problema che nessuna singola entità può risolvere da sola.

\section{Metodologia}
Stessa formula del Lab~1 (scale complete nel Capitolo~\ref{lab:01}):
\[
SE = DPC \times EI + CB
\]
con la mappatura di lavoro $SE<2$ solo registro, $2 \le SE <3$ CNIL, $SE \ge 3$ CNIL e iscritti
(salvo art.~34(3)). In parallelo corre un \textbf{secondo orologio}: secondo il caso, ai sensi del
DORA il fondo deve inviare all'ACPR una notifica iniziale entro 24 ore, intermedia entro 72 ore e
finale entro un mese. È una questione distinta dalla notifica GDPR.

\section{Svolgimento ragionato}

\begin{integrazione}[Svolgimento dei redattori -- non è una soluzione ufficiale]
\paragraph{Qualificazione.} È un data breach ai sensi dell'art.~4(12): esfiltrazione confermata di
dati personali degli iscritti (riservatezza) da parte di un terzo non autorizzato. Il fatto che il
vettore d'attacco sia nei sistemi del fornitore non toglie al fondo, in quanto titolare, gli obblighi
degli artt.~33 e 34; incide sulla responsabilità del fornitore (art.~28) e sui rapporti contrattuali.

\paragraph{Perimetri.} Perimetro~A: ciò che risulta preso (G1, più G2 per quanto si saprà). Perimetro~B:
tutto ciò che il server riceveva, incluso l'export giornaliero completo (G3), dove dati storici e
recenti non si separano. G4 è escluso (accesso fermo alla DMZ), G5 vuoto allo stato.

\paragraph{Punteggio di G1.}
\begin{itemize}
  \item \textbf{DPC}: dati finanziari, base 3. Il valore della posizione dà un'indicazione concreta
    della situazione patrimoniale, quindi non è un ``semplice numero di conto'' (che giustificherebbe
    2). Abbassarlo a 2 è sostenibile solo argomentando che un singolo importo, senza movimenti, dà un
    quadro limitato. Potrebbe salire a 4 se l'export storico permette di ricostruire l'andamento nel
    tempo (volume).
  \item \textbf{EI}: nome e cognome, nel contesto di un fondo specifico e associati a un importo:
    identificazione significativa o diretta (0,75--1).
  \item \textbf{CB}: riservatezza +0,5 (esfiltrazione verso un soggetto criminale, destinatari non
    delimitabili; +0,25 sostenibile se si ritiene il destinatario unico); intento malevolo +0,5;
    integrità e disponibilità 0 (nessuna alterazione, continuità intatta). CB = 1,0.
\end{itemize}
\begin{center}
\small
\begin{tabular}{cccll}
\toprule
\textbf{DPC} & \textbf{EI} & \textbf{CB} & \textbf{SE} & \textbf{Livello} \\ \midrule
2 & 1 & 1,0 & 3,0 & alto (al limite) \\
3 & 0,75 & 1,0 & 3,25 & alto \\
3 & 1 & 1,0 & 4,0 & molto alto \\
\bottomrule
\end{tabular}
\end{center}
Flag: più di 100 interessati. Il caso è costruito per stare \textbf{sul confine} tra medio e alto.

\paragraph{Severità e probabilità.} L'accesso a G1 è certo; l'abuso non è accertato, ma la
combinazione nome + importo è esattamente ciò che serve per un attacco di social engineering
credibile (``la chiamo dal fondo, la sua posizione è di X euro''). La curveball~1 aumenta la
probabilità: più società finanziarie colpite attraverso lo stesso fornitore suggeriscono un'operazione
organizzata e a scopo di lucro.

\paragraph{Decisione.}
\begin{itemize}
  \item \textbf{CNIL}: notifica \textbf{preliminare} (art.~33(4)), senza attendere il quadro completo,
    con impegno a integrarla. La posizione del CEO va respinta: protezione con certificati e colpa del
    fornitore non rilevano, perché l'attaccante ha usato le credenziali legittime del fornitore e i
    dati sono comunque usciti; ``nessun abuso rilevato'' non equivale a ``rischio improbabile''.
  \item \textbf{Iscritti}: la tesi dell'Head of Member Services è forte. Una comunicazione (art.~34)
    con consigli pratici, ad esempio ``il fondo non vi chiederà mai password o pagamenti al
    telefono'', riduce proprio il rischio principale. In alternativa, chi sceglie ``solo CNIL'' deve
    documentare i trigger che farebbero scattare la comunicazione.
  \item \textbf{ACPR}: notifica DORA separata, già avviata con la classificazione delle 11:12.
  \item \textbf{Fornitore}: richiesta formale di informazioni (categorie, volume, finestra, log),
    contestazione della settimana di silenzio (art.~33(2)), riserva di azioni contrattuali.
  \item \textbf{Registro} (art.~33(5)): tutte le scelte e i relativi motivi.
\end{itemize}

\paragraph{Le 72 ore.} Per il titolare decorrono da quando \emph{il titolare} è venuto a conoscenza:
venerdì 11/4 alle 10:44 (comunicazione del responsabile), quindi scadenza \textbf{lunedì 14/4 alle
10:44} (il weekend conta). Questo regge solo se il fondo può dimostrare di non aver potuto sapere
prima: nessuna visibilità sul server durante la settimana silenziosa, nessun alert. Il ritardo del
fornitore (conoscenza il 4/4) è una sua violazione dell'art.~33(2), non sposta indietro l'orologio
del titolare. È lo schema anticipato dal prof.~Vaciago il 05/10 (Capitolo~\ref{lez:2026-10-05}).

\paragraph{La curveball~1.} La società sorella è un \textbf{titolare distinto} dei dati dei propri
clienti: valuta e notifica per conto proprio, e per i suoi clienti di alto profilo il DPC sale
(caratteristiche particolari degli interessati). Il punteggio di Sylvane, sul proprio perimetro, non
cambia; cambiano la probabilità di abuso e la dimensione di gruppo: comunicazioni coerenti (se la
sorella informa i clienti e Sylvane no, il messaggio è incoerente), un'unica linea verso il fornitore
condiviso, coordinamento della risposta.
\end{integrazione}

\section{Indicazioni emerse nel debrief (19/10)}
\begin{esame}
\begin{itemize}
  \item \textbf{Riservatezza senza esfiltrazione}: la violazione di riservatezza c'è anche se i file
    sono stati solo visti e non portati via; l'impatto sugli interessati, però, è molto diverso
    (prof.~Vaciago).
  \item \textbf{Le fasce di rischio sono una guida, non una regola matematica}: nel worksheet le fasce
    ``$2 \le 3$'' e ``$3 \le 4$'' si sovrappongono proprio su 3. Con $SE = 3$ si è in zona grigia: si
    decide e ci si assume la responsabilità in base all'\emph{accountability}; l'autorità può valutare
    diversamente. Gruppi diversi, sullo stesso caso, sono arrivati a numeri molto diversi.
  \item \textbf{Non presumere nulla che non si possa provare} (Gianluca): ogni dichiarazione falsa
    all'autorità è un reato e chi la firma ne risponde personalmente. Nel dubbio ci si attiene alle
    evidenze. Il prof.~Vaciago aggiunge che serve il dolo (\emph{mens rea}): l'errore in buona fede
    non è punito, ma va poi spiegato a un pubblico ministero spesso digiuno di tecnica.
  \item \textbf{Responsabilità del fornitore}: se il vettore d'attacco è nell'infrastruttura del
    responsabile, è sua la colpa di essere stato violato, anche se i dati stavano sul server del
    titolare.
\end{itemize}
\end{esame}

\paragraph{Una delle soluzioni presentate.} Un gruppo ha diviso i dati in due insiemi (canale di
scambio; una copia di backup sul server con dati più ricchi, senza evidenza di accesso nei log),
notificando CNIL e iscritti (lettera con l'avvertimento che il fondo non chiederà mai password o
denaro) e segnalando all'ACPR. Ha fatto partire le 72 ore dall'11/4, potendo dimostrare di non aver
saputo prima, e ha proposto una risposta condivisa a livello di gruppo, documentando e firmando ogni
passaggio. Il docente ha giudicato la presentazione molto buona.

\section{Errori tipici}
\begin{itemize}
  \item Accettare l'argomento ``è colpa del fornitore'' come motivo per non notificare: gli obblighi
    di notifica restano del titolare.
  \item Far partire le 72 ore dal 4/4 (conoscenza del fornitore) oppure, all'opposto, rimandarle in
    attesa di dettagli: la notifica preliminare esiste proprio per questo.
  \item Trattare la fascia di rischio come una regola matematica, specialmente con $SE = 3$.
  \item Fondare la decisione su ipotesi non documentate (es.\ ``probabilmente sono solo dati
    recenti'').
  \item Confondere la notifica GDPR alla CNIL con la segnalazione DORA all'ACPR: due binari, due
    autorità, due orologi.
  \item Valutare voce per voce invece di dare un giudizio complessivo sulla violazione.
  \item Ignorare la curveball o, all'opposto, assorbire nel proprio perimetro i dati della società
    sorella, che ha un proprio titolare.
\end{itemize}

\section{Collegamenti}
\begin{itemize}
  \item GDPR: artt.~4(12), 28, 33(2), 33(4), 33(5), 34 (Capitolo~\ref{cap:data-breach}).
  \item DORA: notifica degli incidenti gravi all'autorità finanziaria (Capitolo~\ref{cap:dora}).
  \item Lezione del 05/10: ritardo del responsabile e decorrenza delle 72 ore; art.~168 Codice privacy
    sulle false dichiarazioni (Capitolo~\ref{lez:2026-10-05}).
  \item Lab~1: metodologia e scale ENISA (Capitolo~\ref{lab:01}).
\end{itemize}

\begin{verifica}
\begin{itemize}
  \item \textbf{Curveball successive}: nel materiale c'è solo la curveball~1. Dal debrief del 19/10
    (trascrizione molto disturbata) sembrano esserne emerse altre: un account non autorizzato presente
    da tempo, il coinvolgimento di altre istituzioni finanziarie e una \textbf{copia di backup} sul
    server con dati più ricchi (identificativi fiscali, dati KYC, transazioni) senza evidenza di
    accesso. Da ricostruire con certezza.
  \item \textbf{Scadenze DORA}: il caso indica notifica iniziale entro 24 ore dalla rilevazione,
    intermedia entro 72 ore, finale entro un mese. Le norme tecniche attuative del DORA fissano la
    notifica iniziale a 4 ore dalla classificazione (e comunque entro 24 ore dalla conoscenza):
    verificare nel capitolo DORA quale formulazione usa il corso.
  \item La data dell'esercizio è l'11 aprile \textbf{2025} (come nel materiale), non 2026.
  \item Lo svolgimento ragionato non è stato validato dai docenti: i punteggi sono alternative
    argomentabili, non la ``risposta giusta''.
\end{itemize}
\end{verifica}
